% Generated draft: requires mathematical and source review.
\documentclass[11pt]{article}
\usepackage[margin=25mm]{geometry}
\usepackage{fontspec}
\setmainfont{Latin Modern Roman}
\usepackage{amsmath,amssymb,mathtools,unicode-math}
\setmathfont{Latin Modern Math}
\usepackage{enumitem,needspace,xurl,hyperref}
\setlength{\parindent}{0pt}
\setlength{\parskip}{5pt}
\setlength{\emergencystretch}{3em}
\newcommand{\N}{\mathbb{N}}
\newcommand{\Z}{\mathbb{Z}}
\newcommand{\Q}{\mathbb{Q}}
\newcommand{\R}{\mathbb{R}}
\newcommand{\C}{\mathbb{C}}
\newcommand{\F}{\mathbb{F}}
\newcommand{\A}{\mathbb{A}}
\newcommand{\PP}{\mathbb{P}}
\newcommand{\E}{\mathbb{E}}
\newcommand{\eps}{\varepsilon}
\newcommand{\dd}{\,\mathrm{d}}
\begin{document}
\section*{UnsolvedMath ID: 800230}
\textbf{Source problem code:} catalogue entry
\subsection{Definitions and mathematical statement}
Let \(\Omega\subset\C\) be a bounded domain with smooth boundary, and use planar area measure \(\dd A\). Let \(L^2(\Omega)\) denote the complex Hilbert space with inner product \(\langle u,v\rangle=\int_\Omega u\overline v\,\dd A\), and let
\[
A^2(\Omega)=\{g\in L^2(\Omega):g\text{ is holomorphic on }\Omega\}.
\]
Interpret the Kohn solution operator \(T:L^2(\Omega)\to L^2(\Omega)\) as the canonical solution operator characterized by
\[
\frac{\partial(Tf)}{\partial\bar z}=f\quad\text{in distributions},\qquad Tf\perp A^2(\Omega).
\]
For \(h\in L^\infty(\Omega)\), define
\[
S_h f=\overline h\,Tf.
\]
The source describes \(S_h\) as compact and non-self-adjoint; self-adjoint exceptional cases, such as \(h\equiv0\), are possible.

The general question is to determine conditions on \(\Omega\) and \(h\) under which \(S_h\) has eigenvalues. For the principal question, interpret this as asking for a \emph{nonzero} eigenvalue, namely whether
\[
\exists\lambda\in\C\setminus\{0\}\ \exists f\in L^2(\Omega)\setminus\{0\}:\quad \overline h\,Tf=\lambda f.
\]
This interpretation is necessary to qualify the source's unqualified use of ``eigenvalues'': when \(h\equiv0\), zero is an eigenvalue even though \(h\) has a single-valued holomorphic primitive.

In particular, suppose \(h\) is bounded and holomorphic on \(\Omega\). A single-valued holomorphic primitive is a holomorphic function \(H:\Omega\to\C\) satisfying \(H'=h\). The source states that the existence of such an \(H\) excludes nonzero eigenvalues of \(S_h\), and asks whether the converse holds. Equivalently, is the following implication valid for every such \(\Omega\) and \(h\)?
\[
\left[\exists\text{ a closed piecewise }C^1\text{ curve }\gamma\subset\Omega:
\int_\gamma h(z)\,\dd z\ne0\right]
\quad\Longrightarrow\quad
\left[\exists\lambda\in\C\setminus\{0\}\ \exists f\in L^2(\Omega)\setminus\{0\}:S_hf=\lambda f\right].
\]
Here the source's phrase ``periods of \(H\)'' is interpreted as the additive monodromy of a locally defined primitive, measured by the periods \(\int_\gamma h(z)\,\dd z\). The source reports an affirmative answer for annuli. That special-case assertion and the primitive obstruction are supplied source claims, not independently verified results in this draft.
\subsection{Short English statement}
For the canonical \(\bar\partial\)-solution operator on a smooth bounded planar domain, determine when multiplication of its output by \(\overline h\) produces a nonzero eigenvalue. When \(h\) is bounded and holomorphic, must a nonzero period of \(h(z)\,\dd z\), equivalently the absence of a single-valued holomorphic primitive, force such an eigenvalue? The source reports that this holds for annuli. The qualification ``nonzero'' resolves an ambiguity in the supplied wording.
\subsection{Sources}
\begin{itemize}
\item Local UnsolvedMath registry: supplied complete record, permanent ID \(800230\), exact source problem code \texttt{catalogue entry}, title ``Several Complex Variables Problem List --- Problem 27.'' The registry supplies the status \texttt{open}; this label is not an independently verified current mathematical assessment.
\item Buzzard, Fornaess, and Krantz, \emph{The Several Complex Variables Problem List} (1995), Problem 27. Supplied source URL: \url{https://arxiv.org/abs/math/9509202}. The registry records extraction by \texttt{arxiv\_latex} and access on 2026-10-16. This access date is supplied provenance and does not indicate independent source access or verification for this draft.
\item The problem text attributes the question to D. Barrett, with the supplied contact \texttt{barrett@math.lsa.umich.edu}. The registry's \texttt{proposed\_by} field instead names the list authors, Buzzard, Fornaess, and Krantz; both attribution levels are retained here.
\end{itemize}
\end{document}
